% !TEX root = ../proyect.tex

\chapter{Conclusiones}\label{cap:conclusiones}

\section{Objetivos Alcanzados}\label{sec:obj-alcanzados}

El desarrollo de este proyecto ha permitido alcanzar los objetivos planteados en el cap\'itulo~\ref{cap:introduccion}:

Se ha dise\~nado e implementado una base de datos en MongoDB con siete colecciones que modelan todas las entidades del dominio: usuarios, profesionales, servicios, citas, bloqueos temporales, notas t\'ecnicas y configuraci\'on global. Los esquemas de Mongoose garantizan la validaci\'on de datos, las restricciones de tipo y formato, y la integridad referencial a nivel de aplicaci\'on.

Se ha desarrollado un portal p\'ublico de reservas funcional que permite a los clientes consultar el cat\'alogo de servicios, visualizar la disponibilidad en tiempo real y realizar reservas de forma aut\'onoma mediante un flujo guiado de tres pasos. La interfaz se adapta a dispositivos m\'oviles gracias al dise\~no \textit{responsive} con Tailwind CSS.

Se ha implementado un panel de administraci\'on completo que proporciona a los gestores del negocio las herramientas necesarias para el control de la agenda, incluyendo un calendario interactivo con vista por profesional basado en FullCalendar, fichas de cliente con notas t\'ecnicas categorizadas (CRM), gesti\'on de horarios semanales con descansos, administraci\'on de bloqueos temporales, y un \textit{dashboard} con indicadores de rendimiento.

El motor de disponibilidad calcula los huecos libres considerando horarios de cada profesional, descansos, bloqueos activos y citas existentes, garantizando la imposibilidad de solapamiento de citas. Este c\'alculo tiene en cuenta la zona horaria configurable del negocio y la rejilla de \textit{slots} definida en los ajustes.

\section{Conocimientos Aplicados}\label{sec:conocimientos}

Este proyecto ha permitido aplicar de forma pr\'actica los conocimientos adquiridos en la asignatura de Complementos de Bases de Datos, especialmente en las siguientes \'areas:

\begin{itemize}
    \item Modelado de datos orientado a documentos con MongoDB y Mongoose.
    \item Dise\~no de esquemas con subdocumentos embebidos, referencias entre colecciones y propiedades virtuales.
    \item Definici\'on de \'indices compuestos para la optimizaci\'on de consultas frecuentes.
    \item Validaci\'on de datos a nivel de esquema y a nivel de aplicaci\'on.
    \item Implementaci\'on de patrones como el \textit{singleton} para la configuraci\'on global con cach\'e en memoria.
    \item Gesti\'on de la integridad referencial en un entorno NoSQL.
\end{itemize}

\section{Dificultades Encontradas}\label{sec:dificultades}

Durante el desarrollo se han encontrado dificultades que han requerido especial atenci\'on:

La gesti\'on de zonas horarias ha supuesto el mayor reto t\'ecnico del proyecto. Todas las fechas se almacenan en UTC en MongoDB, pero los c\'alculos de disponibilidad y la presentaci\'on al usuario deben realizarse en la zona horaria local del negocio. Se han desarrollado utilidades espec\'ificas con Day.js para garantizar la coherencia en las conversiones, evitando errores sutiles como la asignaci\'on de una cita al d\'ia incorrecto en los l\'imites de la medianoche.

El c\'alculo de disponibilidad con duraciones variables ha requerido un dise\~no cuidadoso. Un servicio de 20 minutos en una rejilla de 30 minutos debe bloquear el \textit{slot} completo para evitar fragmentaci\'on de la agenda. La distinci\'on entre duraci\'on real y duraci\'on operativa resuelve este problema.

La sincronizaci\'on de la cach\'e del \textit{frontend} con los cambios en el servidor ha exigido implementar un sistema de invalidaci\'on por grupos mediante React Query, garantizando que todos los componentes reflejen datos actualizados tras cada mutaci\'on.

\section{L\'ineas de Trabajo Futuro}\label{sec:trabajo-futuro}

El sistema desarrollado sienta las bases para futuras ampliaciones:

\begin{itemize}
    \item Implementaci\'on de notificaciones por correo electr\'onico y SMS para recordar citas a los clientes, utilizando servicios como SendGrid o Twilio.
    \item Integraci\'on con pasarelas de pago (Stripe, Redsys) para permitir el cobro anticipado o la se\~nal de reservas.
    \item Desarrollo de un m\'odulo de estad\'isticas avanzadas con gr\'aficos de tendencias de facturaci\'on, servicios m\'as demandados y horarios de mayor afluencia.
    \item Implementaci\'on de bloqueos recurrentes autom\'aticos (el campo \texttt{esRecurrente} ya est\'a contemplado en el modelo).
\end{itemize}

\section{Valoraci\'on Personal}\label{sec:valoracion}

Este proyecto ha supuesto una experiencia enriquecedora, ya que nos ha permitido abordar un problema real y cercano, desarrollando una soluci\'on completa desde el an\'alisis de requisitos hasta la implementaci\'on. El hecho de que el sistema est\'e destinado a ser utilizado en un negocio familiar ha a\~nadido una motivaci\'on adicional y ha permitido obtener retroalimentaci\'on directa de los usuarios finales durante el desarrollo, validando decisiones de dise\~no y ajustando funcionalidades a las necesidades reales del d\'ia a d\'ia del negocio.

